\section*{Chapter \ref{ch:b2:series} --- Sequences and Series}

\begin{solution}{exo:b2:series:1}
$\sum\frac{\cos n}{n}$: Abel's test with $a_n = \frac1n$ and $b_n =
\cos n$, whose partial sums are bounded (real part of a geometric
sum, as in \cref{ex:b2:series:sinn}): convergent (not absolutely,
by the same $\cos^2$ trick).

$\sum \frac{(-1)^n}{\ln n}$ ($n \geq 2$): alternating test,
$\frac{1}{\ln n} \downarrow 0$: convergent; not absolutely ($\ln n
\leq n$).

$\sum \frac{(-1)^n}{n^{3/4} + \cos n}$: expand,
\[
\frac{(-1)^n}{n^{3/4} + \cos n}
= \frac{(-1)^n}{n^{3/4}}\cdot
\frac{1}{1 + \frac{\cos n}{n^{3/4}}}
= \frac{(-1)^n}{n^{3/4}}
- \frac{(-1)^n\cos n}{n^{3/2}} + O\Bigl(\frac{1}{n^{9/4}}\Bigr).
\]
First series: alternating, convergent. Second: absolutely
convergent ($\frac{1}{n^{3/2}}$ scale). Third: absolutely
convergent. Total: convergent.
\end{solution}

\begin{solution}{exo:b2:series:2}
Cauchy product of $\sum_{m\geq0} z^m$ with itself (both absolutely
convergent for $\abs z < 1$): the diagonal coefficient is $c_k =
\sum_{m=0}^{k} 1 = k + 1$, so
\[
\frac{1}{(1-z)^2} = \sum_{k\geq0} (k+1)z^k .
\]
Multiplying by $z$ and reindexing: $\sum_{n \geq 1} n z^n =
\frac{z}{(1-z)^2}$.
\end{solution}

\begin{solution}{exo:b2:series:3}
Let $B_n = \sum_{k \leq n} b_k \to B$. Abel summation with $a_k = k$:
\[
\sum_{k=1}^{n} k\,b_k = n B_n - \sum_{k=1}^{n-1} B_k
\quad\Longrightarrow\quad
\frac1n \sum_{k=1}^{n} k b_k = B_n - \frac{1}{n}\sum_{k=1}^{n-1}
B_k .
\]
The Ces\`aro means of the convergent $(B_k)$ tend to its limit $B$
(Year 1 volume), so the right side tends to $B - B = 0$.
\end{solution}

\begin{solution}{exo:b2:series:4}
If $\theta \in \pi\Z$: all terms vanish --- convergent trivially.
Assume $\theta \notin \pi\Z$.

$\alpha > 1$: absolutely convergent (domination by $n^{-\alpha}$).

$0 < \alpha \leq 1$: Abel's test applies ($a_n = n^{-\alpha}
\downarrow 0$; partial sums of $\sin n\theta$ bounded by
$\frac{1}{\abs{\sin(\theta/2)}}$, geometric sum): convergent. Not
absolutely: $\abs{\sin n\theta} \geq \sin^2 n\theta = \frac{1 -
\cos 2n\theta}{2}$, and $\sum \frac{1 - \cos 2n\theta}{2n^\alpha}$
diverges ($\sum n^{-\alpha}$ diverges; $\sum
\frac{\cos 2n\theta}{n^\alpha}$ converges by Abel when $2\theta
\notin 2\pi\Z$; the excluded case $2\theta \in 2\pi\Z$ means
$\theta \in \pi\Z$, already handled). Semi-convergent.
\end{solution}

\begin{solution}{exo:b2:series:5}
Let $p_1, p_2, \dots$ be the nonnegative terms of $(u_n)$ in order,
$q_1, q_2, \dots$ the negative ones. Both $\sum p_k$ and $\sum
q_k$ diverge: if one of them converged, the other would equal the
convergent $\sum u_n$ minus it, hence converge too --- and then
$\sum \abs{u_n} = \sum p_k - \sum q_k$ would converge, contradicting
the hypothesis. Also $u_n \to 0$ ($\sum u_n$ converges).

Greedy rearrangement: take positive terms $p_1, p_2, \dots$ until
the running total first exceeds $\ell$ (possible: $\sum p_k =
+\infty$); then negative terms until the total first drops below
$\ell$ (possible: $\sum q_k = -\infty$); repeat forever (each
phase is finite, and every term is used exactly once: a genuine
rearrangement). After each switch, the distance from the running
total to $\ell$ is at most the last term used; since the terms used
at the $m$-th switch have index $\to \infty$, and $u_n \to 0$,
the running totals converge to $\ell$.
\end{solution}

\begin{solution}{exo:b2:series:6}
Summability: the finite partial sums of
$\frac{\abs x^{m+n}}{m!n!}$ are bounded by $\bigl(\sum_m
\frac{\abs x^m}{m!}\bigr)^2 = \eu^{2\abs x}$. Diagonal grouping
(\cref{thm:b2:series:fubini}):
\[
(\eu^{x})^2 = \sum_{m,n} \frac{x^{m+n}}{m!\,n!}
= \sum_{k=0}^{\infty} x^k \sum_{m+n=k} \frac{1}{m!\,n!}
= \sum_k \frac{x^k}{k!}\sum_{m=0}^{k}\binom km
= \sum_k \frac{(2x)^k}{k!} = \eu^{2x} .
\]
\end{solution}

\begin{solution}{exo:b2:series:7}
Summability: bounded by $\zeta(2)^2$ as a product family
(\cref{thm:b2:series:fubini}'s Cauchy-product argument). The map
$(q, a, b) \mapsto (qa, qb)$, from triples with $\gcd(a, b) = 1$
to pairs $(m, n)$, is a bijection (set $q = \gcd(m,n)$). Grouping
the summable family accordingly (a partition of the index set ---
legitimate for summable families by
\cref{thm:b2:series:rearrangement}/\cref{thm:b2:series:fubini}
applied to the partition into countably many classes):
\[
\zeta(2)^2 = \sum_{q} \sum_{\gcd(a,b)=1} \frac{1}{q^4 a^2 b^2}
= \zeta(4)\, S,
\qquad\text{so}\qquad
S = \frac{\zeta(2)^2}{\zeta(4)} .
\]
(With the values $\zeta(2) = \frac{\pi^2}{6}$,
$\zeta(4) = \frac{\pi^4}{90}$ from
\cref{ch:b2:fourier}: $S = \frac{5}{2}$.)
\end{solution}

\begin{solution}{exo:b2:series:8}
Let $A = \sum a_n$ (absolute), $B_m = \sum_{k\leq m} b_k \to B$,
bounded by $M$. Then
\[
C_N = \sum_{k=0}^{N} c_k = \sum_{n=0}^{N} a_n B_{N-n}
\]
(collect by the index of $a$). Write
\[
C_N - AB = \sum_{n=0}^{N} a_n (B_{N-n} - B) - B\sum_{n > N} a_n .
\]
The last term tends to $0$. Split the sum at $n = \lfloor N/2
\rfloor$: for $n \leq N/2$, $N - n \geq N/2$, so $\abs{B_{N-n} - B}
\leq \varepsilon_N := \sup_{m \geq N/2}\abs{B_m - B} \to 0$, and
this part is $\leq \varepsilon_N \sum\abs{a_n}$; for $n > N/2$,
$\abs{B_{N-n} - B} \leq 2M$, and this part is $\leq 2M \sum_{n >
N/2} \abs{a_n} \to 0$. Hence $C_N \to AB$.
\end{solution}

\begin{solution}{exo:b2:series:9}
Take $u_n = 2^{-n} e_n$ ($e_n$ the sequence with a single $1$ in
position $n$). Then $\sum \norm{u_n}_\infty = \sum 2^{-n} <
\infty$: absolutely convergent. But the partial sums $S_N =
(1, \tfrac12, \dots, 2^{-N}, 0, \dots)$ would have to converge to
the sequence $(2^{-n})_n$, which is \emph{not} eventually zero:
outside $E$. Inside $E$, $(S_N)$ is Cauchy without limit ($\norm{S_N
- x}_\infty \geq 2^{-N-1}$ fails to help any eventually-zero $x$:
for any $x \in E$ vanishing beyond rank $K$, $\norm{S_N - x} \geq
2^{-K-1}$ for $N > K$): the series does not converge in $E$.
Completeness is exactly what \cref{thm:b2:nvs:absoluteconvergence}
needs.
\end{solution}

\begin{solution}{exo:b2:series:10}
Both $\arctan(n+1) - \arctan n$ and
$\arctan\frac{1}{n^2+n+1}$ lie in $\intoo{0}{\frac\pi2}$, and the
tangent addition formula gives
\[
\tan\bigl(\arctan(n{+}1) - \arctan n\bigr)
= \frac{(n+1) - n}{1 + n(n+1)} = \frac{1}{n^2 + n + 1} :
\]
equal tangents in an interval where $\tan$ is injective, so the
identity holds. Telescoping,
\[
\sum_{n=1}^{N}\arctan\frac{1}{n^2+n+1}
= \arctan(N{+}1) - \arctan 1 \xrightarrow[N\to\infty]{}
\frac\pi2 - \frac\pi4 = \frac\pi4 .
\]

\end{solution}

\begin{solution}{exo:b2:series:11}
First: $\sqrt{n+1} - \sqrt n = \frac{1}{\sqrt{n+1} + \sqrt n}
\sim \frac{1}{2\sqrt n}$, so the terms are $\sim
2^{-\alpha}n^{-\alpha/2}$: convergence iff $\frac\alpha2 > 1$,
i.e.\ $\alpha > 2$. Second: $1 - \cos\frac1n \sim
\frac{1}{2n^2}$: converges. Third: $\bigl(1 + \frac1n\bigr)^n =
\eu^{\,n\ln(1 + 1/n)} = \eu^{\,1 - \frac1{2n} + O(n^{-2})} =
\eu\bigl(1 - \frac{1}{2n} + O(n^{-2})\bigr)$, so
\[
\eu - \Bigl(1 + \frac1n\Bigr)^{\!n} \sim \frac{\eu}{2n} :
\]
positive terms equivalent to a harmonic multiple: diverges.
\end{solution}

\begin{solution}{exo:b2:series:12}
By monotonicity, $n\,a_{2n} \leq a_{n+1} + a_{n+2} + \dots +
a_{2n} = S_{2n} - S_n \to 0$ (Cauchy criterion for the
convergent series). Hence $2n\,a_{2n} \to 0$, and $(2n{+}1)\,
a_{2n+1} \leq (2n{+}1)a_{2n} = \frac{2n+1}{2n}\,(2n\,a_{2n}) \to
0$: both subsequences of $(na_n)$ tend to $0$, so $na_n \to 0$.

\emph{Converse fails:} $a_n = \frac1{n\ln n}$ is positive
decreasing with $na_n = \frac1{\ln n} \to 0$, yet $\sum a_n$
diverges (Bertrand frontier, \cref{pb:b2:series:1}, question
18). \emph{Monotonicity essential:} let $a_n = \frac1n$ when $n$
is a power of $2$ and $a_n = 2^{-n}$ otherwise: $\sum a_n \leq
\sum_k 2^{-k} + \sum_n 2^{-n} < \infty$, but $na_n = 1$ along
the powers of $2$: $na_n \not\to 0$.
\end{solution}

\begin{solution}{pb:b2:series:1}
\textbf{1.} Induction on $m$: for $m = 0$ both sides are $1$;
the step multiplies by $\cos\theta + \iu\sin\theta$ and uses the
addition formulas $\cos(m\theta + \theta) =
\cos m\theta\cos\theta - \sin m\theta\sin\theta$, $\sin(m\theta
+ \theta) = \sin m\theta\cos\theta + \cos m\theta\sin\theta$.
Expanding instead by the binomial theorem with $m = 2n+1$ and
collecting the imaginary part (the odd powers of
$\iu\sin\theta$, with $\iu^{2j+1} = (-1)^j\iu$):
\[
\sin\bigl((2n{+}1)\theta\bigr) = \sum_{j=0}^{n}(-1)^j
\binom{2n+1}{2j+1}\cos^{2(n-j)}\theta\,\sin^{2j+1}\theta .
\]

\textbf{2.} On $\intoo0{\frac\pi2}$, $\sin\theta \neq 0$: factor
$\sin^{2n+1}\theta$ from each term, leaving
$\bigl(\frac{\cos^2\theta}{\sin^2\theta}\bigr)^{n-j} =
(\cot^2\theta)^{n-j}$: the displayed identity with $P_n(x) =
\sum_j(-1)^j\binom{2n+1}{2j+1}x^{n-j}$. Its degree-$n$
coefficient is $j = 0$'s $\binom{2n+1}{1} = 2n + 1 \neq 0$.

\textbf{3.} At $\theta_k = \frac{k\pi}{2n+1}$:
$\sin\bigl((2n{+}1)\theta_k\bigr) = \sin k\pi = 0$ while
$\sin^{2n+1}\theta_k \neq 0$, so $P_n(\cot^2\theta_k) = 0$. The
$\theta_k$ increase strictly in $\intoo0{\frac\pi2}$, where
$\cot^2$ is strictly decreasing: the values $x_k =
\cot^2\theta_k$ are pairwise distinct --- $n$ distinct roots of
a degree-$n$ polynomial, hence all of them.

\textbf{4.} Vieta: the sum of the roots is minus the ratio of
the $x^{n-1}$- and $x^n$-coefficients:
\[
\sum_{k=1}^{n}\cot^2\theta_k =
\frac{\binom{2n+1}{3}}{\binom{2n+1}{1}}
= \frac{(2n+1)(2n)(2n-1)/6}{2n+1} = \frac{n(2n-1)}{3}.
\]

\textbf{5.} $\frac{1}{\sin^2\theta} = 1 + \cot^2\theta$:
summing, $n + \frac{n(2n-1)}3 = \frac{3n + 2n^2 - n}{3} =
\frac{2n(n+1)}{3}$.

\textbf{6.} On $\intoo{0}{\frac\pi2}$: $\sin\theta < \theta <
\tan\theta$ (Year 1 volume). Taking reciprocals reverses:
$\cot\theta < \frac1\theta < \frac1{\sin\theta}$, and squaring
(all positive) gives $\cot^2\theta < \frac1{\theta^2} <
\frac1{\sin^2\theta}$.

\textbf{7.} Sum question 6 at $\theta = \theta_k$ over $k \leq
n$, using questions 4 and 5, and $\frac1{\theta_k^2} =
\frac{(2n+1)^2}{k^2\pi^2}$:
\[
\frac{n(2n-1)}3 < \frac{(2n+1)^2}{\pi^2}\sum_{k=1}^n\frac1{k^2}
< \frac{2n(n+1)}3 .
\]

\textbf{8.} Multiply by $\frac{\pi^2}{(2n+1)^2}$:
\[
\frac{\pi^2}{3}\cdot\frac{n(2n-1)}{(2n+1)^2}
< \sum_{k=1}^{n}\frac1{k^2} <
\frac{\pi^2}{3}\cdot\frac{2n(n+1)}{(2n+1)^2}.
\]
Both bounds tend to $\frac{\pi^2}3\cdot\frac12 = \frac{\pi^2}6$
(the rational fractions tend to $\frac12$). The partial sums
increase, so they converge, and the squeeze gives $\zeta(2) =
\frac{\pi^2}6$: Euler's theorem, by Cauchy's proof.

\textbf{9.} The partial sums increase to $\zeta(2) =
\frac{\pi^2}6$, so $0 \leq \frac{\pi^2}6 - \sum_{k\leq n}k^{-2}$;
and the lower bound of question 8 gives
\[
\frac{\pi^2}6 - \sum_{k\leq n}\frac1{k^2}
\leq \frac{\pi^2}6 - \frac{\pi^2}3\cdot\frac{n(2n-1)}{(2n+1)^2}
= \frac{\pi^2}6\cdot\frac{(2n+1)^2 - (4n^2 -
2n)}{(2n+1)^2}
= \frac{\pi^2}6\cdot\frac{6n + 1}{(2n+1)^2}
= O\Bigl(\frac1n\Bigr),
\]
matching the exact tail $\frac1n - \frac1{2n^2} + O(n^{-3})$ of
\cref{exo:b2:comparison:11}.

\textbf{10.} The second elementary symmetric function of the
roots is $\sigma_2 = \frac{\binom{2n+1}5}{\binom{2n+1}1} =
\frac{(2n)(2n-1)(2n-2)(2n-3)}{120}$, so
\[
\sum_k\cot^4\theta_k = \sigma_1^2 - 2\sigma_2
= \Bigl(\frac{n(2n-1)}3\Bigr)^2 -
\frac{2n(2n-1)(2n-2)(2n-3)}{60}
\sim \frac{4n^4}9 - \frac{4n^4}{15} = \frac{8n^4}{45}.
\]
Squeezing $\cot^4\theta < \theta^{-4} < (1 + \cot^2\theta)^2 =
1 + 2\cot^2\theta + \cot^4\theta$ and summing: both outer sums
are $\frac{8n^4}{45}(1 + o(1))$ (the added $n + 2\sigma_1 =
O(n^2)$ is negligible), while the middle is
$\frac{(2n+1)^4}{\pi^4}\sum_{k\leq n}k^{-4}$. Hence
\[
\sum_{k\leq n}\frac1{k^4} \longrightarrow
\pi^4\cdot\frac{8/45}{16} = \frac{\pi^4}{90}.
\]

\textbf{11.} Splitting $\zeta(2)$ over parities: $\sum_{\text{even}}
= \sum_j\frac{1}{(2j)^2} = \frac14\zeta(2) = \frac{\pi^2}{24}$,
so $\sum_{\text{odd}} = \zeta(2) - \frac{\pi^2}{24} =
\frac{\pi^2}8$. Alternating: $\sum_k\frac{(-1)^{k-1}}{k^2} =
\sum_{\text{odd}} - \sum_{\text{even}} = \frac{\pi^2}8 -
\frac{\pi^2}{24} = \frac{\pi^2}{12}$ (absolute convergence
justifies the regrouping,
\cref{thm:b2:series:rearrangement}).

\textbf{12.} $S = \frac{\zeta(2)^2}{\zeta(4)} =
\frac{(\pi^2/6)^2}{\pi^4/90} = \frac{90}{36} = \frac52$.
Heuristic: the identity $\zeta(2)^2 = \zeta(4)S$ of
\cref{exo:b2:series:7} says that pulling out the $\gcd$
renormalizes pairs into coprime pairs; the reciprocal
$\frac1{\zeta(2)} = \frac6{\pi^2} \approx 0.608$ is the natural
candidate for the density of coprime pairs among all pairs ---
a statement about $\lim_N \frac{1}{N^2}\#\{(m,n) \leq N :
\gcd = 1\}$ whose honest proof (with error terms) belongs to
the Year 3 volume.

\textbf{13.} Direct summation has error $\sim \frac1n$: six
digits require about $10^6$ terms. The corrected sum
$\sum_{k\leq n}k^{-2} + \frac1n - \frac1{2n^2}$ has error
$O(n^{-3})$: at $n = 100$ it equals $1.6449339\dots$ against
$\frac{\pi^2}6 = 1.6449341\dots$ --- error $1.7\cdot10^{-7}$,
seven digits from a hundred terms. Asymptotic corrections beat
raw patience by four orders of magnitude.

\textbf{14.} $n = 1$: $P_1(x) = 3x - 1$, root $\frac13$, and
indeed $\cot^2\frac\pi3 = \bigl(\frac1{\sqrt3}\bigr)^2 =
\frac13 = \frac{1\cdot1}3$. $n = 2$: the formula predicts
$\frac{2\cdot3}3 = 2$; with $\cos\frac\pi5 = \frac{1 +
\sqrt5}{4}$, one computes $\cot^2 36^\circ \approx 1.894$ and
$\cot^2 72^\circ \approx 0.106$: sum $2.000$.

\textbf{15.} The numbers $\tan^2\theta_k = \frac1{x_k}$ are the
roots of $Q(x) = x^nP_n\bigl(\frac1x\bigr) =
\sum_{j=0}^n(-1)^j\binom{2n+1}{2j+1}x^j$ (the $x_k$ are
nonzero). Vieta on $Q$: the leading coefficient is $(-1)^n$
(term $j = n$), the next is $(-1)^{n-1}\binom{2n+1}{2n-1} =
(-1)^{n-1}\binom{2n+1}{2}$, so
\[
\sum_{k=1}^n\tan^2\theta_k =
-\frac{(-1)^{n-1}\binom{2n+1}2}{(-1)^n} = \binom{2n+1}2\cdot
\frac{2}{2n+1}\cdot\frac{2n+1}{2}
= n(2n+1).
\]
Check $n = 1$: $\tan^2\frac\pi3 = 3 = 1\cdot3$.

\textbf{16.} Let $\gamma = \frac{1+\alpha}2 \in
\intoo{1}{\alpha}$. Then $\frac{n^{-\alpha}(\ln
n)^{-\beta}}{n^{-\gamma}} = n^{\gamma - \alpha}(\ln n)^{-\beta}
\to 0$ (a negative power of $n$ beats any power of $\ln n$), so
eventually the terms are $\leq n^{-\gamma}$ with $\gamma > 1$:
convergence by comparison with a Riemann series.

\textbf{17.} Let $\gamma = \frac{1+\alpha}2 \in
\intoo{\alpha}{1}$: now $\frac{n^{-\gamma}}{n^{-\alpha}(\ln
n)^{-\beta}} = n^{\alpha-\gamma}(\ln n)^{\beta} \to 0$, so
eventually the terms are $\geq n^{-\gamma}$ with $\gamma < 1$:
divergence.

\textbf{18.} $f(t) = \frac1{t(\ln t)^\beta}$ is positive,
continuous, and decreasing for large $t$ (its logarithm has
derivative $-\frac1t\bigl(1 + \frac{\beta}{\ln t}\bigr) < 0$
eventually). Antiderivatives: for $\beta \neq 1$,
$\int^x f = \frac{(\ln x)^{1-\beta}}{1-\beta} + \text{const}$,
which has a finite limit iff $\beta > 1$; for $\beta = 1$,
$\int^x f = \ln\ln x \to \infty$. By
\cref{thm:b2:comparison:seriesintegral}, the series and the
integral share their nature: convergence iff $\beta > 1$.

\textbf{19.} Same test: $\frac{\dd}{\dd t}\ln\ln\ln t =
\frac{1}{t\ln t\,\ln\ln t}$, and $\ln\ln\ln t \to \infty$:
divergence. And $\frac{\dd}{\dd t}\Bigl(-\frac1{\ln\ln t}\Bigr)
= \frac{1}{t\ln t\,(\ln\ln t)^2}$ with $-\frac1{\ln\ln t} \to
0$: convergence.

\textbf{20.} First: $n^{1/\ln n} = \eu^{\ln n/\ln n} = \eu$, so
the terms are exactly $\frac{1}{\eu\,n}$: a multiple of the
harmonic series, \emph{divergent} --- the exponent $1 +
\frac1{\ln n}$ crawls to $1$ too fast. Second: $n^{1/\ln\ln n}
= \eu^{\ln n/\ln\ln n}$, and $\frac{\ln n}{\ln\ln n} \geq
2\ln\ln n$ eventually, so $n^{1/\ln\ln n} \geq (\ln n)^2$: the
terms are $\leq \frac1{n(\ln n)^2}$, a convergent Bertrand
series (question 18): \emph{convergent}. The frontier passes
strictly between these two exponents.

\textbf{21.} $R_n \downarrow 0$ and
\[
\sqrt{R_n} - \sqrt{R_{n+1}} = \frac{R_n -
R_{n+1}}{\sqrt{R_n} + \sqrt{R_{n+1}}} =
\frac{a_n}{\sqrt{R_n} + \sqrt{R_{n+1}}} \geq
\frac{a_n}{2\sqrt{R_n}},
\]
so $\sum_n \frac{a_n}{\sqrt{R_n}} \leq 2\sum_n(\sqrt{R_n} -
\sqrt{R_{n+1}}) = 2\sqrt{R_1} < \infty$ (telescoping). Yet
$\frac{a_n/\sqrt{R_n}}{a_n} = \frac1{\sqrt{R_n}} \to \infty$:
the new series converges while being infinitely larger. No
convergent series is slowest; comparison tests against any
fixed family can never be complete.

\textbf{22.} For decreasing positive $(a_n)$, group the terms
between consecutive powers of $2$:
\[
2^{k}a_{2^{k+1}} \leq \sum_{n=2^k}^{2^{k+1}-1} a_n \leq
2^ka_{2^k} .
\]
Summing over $k$: if $\sum 2^ka_{2^k}$ converges, the partial
sums of $\sum a_n$ are bounded (converges); if $\sum a_n$
converges, then $\sum_k 2^{k+1}a_{2^{k+1}} \leq 2\sum_n a_n <
\infty$. For $a_n = \frac1{n(\ln n)^\beta}$: $2^ka_{2^k} =
\frac{1}{(k\ln 2)^\beta}$, and $\sum k^{-\beta}$ converges iff
$\beta > 1$: the frontier of question 18 again, without
integrals.

\textbf{23.} By the integral comparison,
\[
\sum_{n > N}\frac{1}{n(\ln n)^2} \geq
\int_{N+1}^{\infty}\frac{\dd t}{t(\ln t)^2} =
\frac{1}{\ln(N+1)},
\]
which at $N = 10^6$ is $\approx 0.0724$: after a million terms
the tail still exceeds $0.07$ --- the series converges, but no
direct summation will ever exhibit its sum. Contrast with
question 13, where one asymptotic correction bought seven
digits from a hundred terms: knowing \emph{how} a series
converges is worth more than knowing that it does.

\textbf{24.} $\sum\frac1{n\ln n}$: diverges ($\alpha = 1$,
$\beta = 1$, question 18). $\sum\frac1{n^{1.01}}$: converges
(Riemann, $\alpha > 1$). $\sum\frac{(\ln n)^{100}}{n^{1.001}}$:
converges ($\alpha = 1.001 > 1$, $\beta = -100$, question 16).
$\sum\frac1{n\ln n(\ln\ln n)^3}$: converges (question 19's
pattern: antiderivative $-\frac12(\ln\ln t)^{-2}$, finite
limit).

\textbf{25.} De Moivre converts the vanishing of $\sin(2n{+}1)
\theta_k$ into the vanishing of a polynomial at
$\cot^2\theta_k$, and Vieta reads off the exact root sums that
analysis alone could only estimate (questions 1--5). The
squeeze needed the \emph{exact} values $\frac{n(2n-1)}3$ and
$\frac{2n(n+1)}3$ on both sides --- equivalents would have
begged the question, since the whole point is the constant
$\frac{\pi^2}6$ (questions 7--8). Part IV ran entirely on
\cref{ch:b2:comparison}'s series--integral comparison, the
logarithmic antiderivatives doing the classifying (questions
18--19). Question 21 destroys the dream of a universal
comparison test: below every convergent series lies another,
infinitely slower --- scales like Bertrand's map the frontier
ever more finely but never reach it. Summits: Euler's $\zeta(2)
= \frac{\pi^2}6$ with its upper floor $\zeta(4) =
\frac{\pi^4}{90}$ (questions 8, 10), and the Bertrand
classification (questions 16--18); $\zeta(2)$ returns in
\cref{ch:b2:fourier}, where Parseval's identity re-proves it
in one line from the Fourier series of the sawtooth --- one
constant, two civilizations.
\end{solution}
